%% ============================================================================
%% ch3_literature.tex -- Chapter 3: סקירת הספרות ומאמר הדגל
%% budget: body pages 8-16 (9 pages). Thematic & critical, NOT paper-by-paper.
%% ============================================================================
\chapter{סקירת הספרות ומאמר הדגל}
\label{ch:lit}

\begin{managerbox}
מי שאינו קורא מה נעשה לפניו משלם על המצאת הגלגל. הפרק מראה שכל אחד מרכיבי \en{AURAM} נשען על שיטה שנבחנה בקהילה המדעית, ומאתר במדויק את הנקודה שבה הידע הקיים נעצר. הנקודה הזו היא פער הידע, והיא ההצדקה לפרויקט כולו.
\end{managerbox}

הסקירה בנויה באופן תמטי וביקורתי \cite{webster2002literature}: במקום לסכם מאמרים לפי סדר כרונולוגי, היא מחלקת את הספרות לארבעה תחומי ידע, ובכל תחום מנתחת את ההנחות המפשטות שעליהן נשענים הפתרונות. איור~\ref{fig:pillars} מציג את ארבעת התחומים ואת נקודות המפגש ביניהם, שבהן נמצא הפרויקט.

\section{שיטת הסקירה ותחומי הידע}
\label{sec:lit-method}

המקורות נבחרו מכנסים וכתבי עת שפיטים מובילים בראייה ממוחשבת, בלמידת מכונה ובגרפיקה (\en{CVPR}, \en{ICCV}, \en{ECCV}, \en{NeurIPS}, \en{ICML}, \en{ACM MM}, \en{ACM TOG}, \en{IEEE TIP}), בדגש על השנים \num{2019--2024}; עבודות יסוד מוקדמות נכללו רק כשהן מגדירות מדד או מודל שהפרויקט משתמש בו ישירות, כמו \en{CIEDE2000} \cite{sharma2005ciede2000}. קריטריון ההכללה היה תרומה ישירה לאחד מארבעת התחומים עם מדדי ביצוע כמותיים, וכל מקור נבדק בשלושת הצעדים של מתודולוגיית הקורס \cite{segal2026booklet}.

\begin{reportfigure}{ארבעת תחומי הידע של הסקירה: הפרויקט ממוקם בחיתוך של כולם}{fig:pillars}
\begin{tikzpicture}[x=1cm,y=1cm,
  p/.style={draw=#1,fill=#1!12,rounded corners=4pt,line width=1pt,minimum width=4.6cm,minimum height=1.35cm,align=center,font=\small}]
  \node[p=brandprimary] (a) at (-3.2,1.25) {\heb{א. קביעות צבע ועוגנים ביומטריים}\\\scriptsize\enlong{Color Constancy, CIEDE2000, Eye Priors}};
  \node[p=brandsecondary] (b) at (3.2,1.25) {\heb{ב. שחזור גוף ויד תלת-ממדי}\\\scriptsize\enlong{SMPL-X, HMR 2.0, Body Measurement}};
  \node[p=brandaccent] (c) at (-3.2,-1.25) {\heb{ג. הדמיה גנרטיבית}\\\scriptsize\enlong{Diffusion VTON, ControlNet, Makeup}};
  \node[p=brandexample] (d) at (3.2,-1.25) {\heb{ד. אחזור רב-מודאלי וטרנדים}\\\scriptsize\enlong{CLIP, RAG, Trend Forecasting}};
  \node[draw=brandwarn,fill=white,circle,line width=1.4pt,minimum size=1.7cm,align=center,font=\scriptsize\bfseries] (x) at (0,0) {\en{AURAM}};
  \foreach \n in {a,b,c,d} \draw[brandwarn,thick,dashed] (x) -- (\n);
\end{tikzpicture}
\end{reportfigure}

כל תחום לבדו מפותח, אך החיבורים ביניהם כמעט אינם נחקרים, כפי שתכמת מטריצת הספרות בסעיף~\ref{sec:matrix}.

\section{תחום א: קביעות צבע ועוגנים ביומטריים}
\label{sec:lit-color}

בעיית קביעות הצבע (\en{Color Constancy}) היא אמידת צבע מקור האור מתוך תמונה, כדי לשחזר את צבע המשטח "האמיתי". המודל הקלאסי הוא מודל פון קריס (\en{von Kries}), שבו כל ערוץ צבע מוכפל במקדם משלו:
\begin{equation}
\mathbf{c}_{\text{obs}} = \operatorname{diag}(\boldsymbol{\ell})\,\mathbf{c}_{\text{true}} + \boldsymbol{\epsilon},
\qquad \boldsymbol{\ell} \in \mathbb{R}^{3}_{+},
\label{eq:vonkries}
\end{equation}
כאשר $\boldsymbol{\ell}$ הוא וקטור צבע המאיר ו-$\boldsymbol{\epsilon}$ רעש החיישן. המודל מניח שהמאיר משנה כל ערוץ בנפרד; ההנחה נשברת במקורות צרי-פס כמו פלורסנט ו-\en{LED} זול, שבהם הספקטרום כולו קובע את התגובה (משוואה~\eqref{eq:formation}). הגישה הלומדת אומדת את המאיר מהתפלגות הצבעים בתמונה: \en{FFCC} בהיסטוגרמה לוגריתמית-כרומטית \cite{barron2017ffcc}, ו-\en{FC4} בשקלול אזורים לפי אמינותם \cite{hu2017fc4}. הגישה המתקנת לאחר מעשה הראתה שכשל של איזון הלבן עובר עיבוד לא-לינארי ואינו מתוקן במכפלה אלכסונית \cite{afifi2019wb}. בצד המדידה, רגרסיה בפולינום-שורש ממפה צבע במדויק, אך דורשת לוח צבעים בתמונה \cite{finlayson2015rootpoly}.

מדד ההצלחה בכל העבודות הללו הוא מרחק בין צבעים במרחב \en{CIELAB}. הנוסחה המקובלת, \en{CIEDE2000}, מתקנת את אי-האחידות התפיסתית של המרחב בעזרת פונקציות משקל לבהירות, לרוויה ולגוון \cite{sharma2005ciede2000}:
\begin{equation}
\Delta E_{00} = \sqrt{\Big(\frac{\Delta L'}{k_L S_L}\Big)^2 + \Big(\frac{\Delta C'}{k_C S_C}\Big)^2 + \Big(\frac{\Delta H'}{k_H S_H}\Big)^2 + R_T\,\frac{\Delta C'}{k_C S_C}\,\frac{\Delta H'}{k_H S_H}} ,
\label{eq:de00}
\end{equation}
כאשר $\Delta L'$, $\Delta C'$ ו-$\Delta H'$ הם הפרשי הבהירות, הרוויה והגוון, $S_L$, $S_C$ ו-$S_H$ פונקציות משקל, ו-$R_T$ איבר סיבוב שמתקן את האזור הכחול. לכן מדד ממוצע אחד עלול להסתיר הטיה לפי גוון עור, וניסוי האימות מנתח תת-קבוצות.

\textbf{ניתוח ביקורתי.} לכל השיטות הללו הנחה משותפת: ההתפלגות הסטטיסטית של צבעי הסצנה מייצגת את העולם שעליו אומן המודל, או שקיים בתמונה ייחוס חיצוני. בתמונת סלפי שתי ההנחות נשברות. העור תופס את רוב הפריים, ולכן אלגוריתם שמניח "עולם אפור" מנטרל דווקא את המידע שמבקשים למדוד. הסימולציה בפרק~\ref{ch:problem} הראתה את זה במספרים: השגיאה עלתה מ-\num{6.58} ל-\num{16.03}. מנגד, אי אפשר לבקש מכל משתמשת להחזיק לוח צבעים.

\textbf{עוגנים ביומטריים.} נישינו ונאיאר הראו כבר ב-\num{2004} שהעין היא "גשש אור" טבעי: מהשתקפות הסביבה בקרנית אפשר לשחזר את תאורת הסצנה \cite{nishino2004eyes}. לובן העין (\en{Sclera}) קיים בכל סלפי, אך הוא אינו כרטיס לבן: הוא מכוסה באפיסקלרה ובלחמית העשירות בכלי דם, ולכן אדמומיות תת-קלינית משנה אותו מיום ליום; עם הגיל שקיעה של שומנים ובילירובין מצהיבה אותו; ושכבת הדמעות מחזירה הבזקים ספקולריים שאינם מפזרים את האור באופן אחיד (לא-למברטיאניים). לכן העוגן חייב להיות אישי ומסונן, והקשתית מספקת עוגן כרומטי שני שהמשתמשת יודעת לתאר. מנגד, תונג, יוניאק ושיאנג הראו שתיאור גוון עור בציר בהירות יחיד אינו מספיק, ושנדרש גם ציר של גוון (\en{Hue}) כדי לייצג הוגנות בין קבוצות \cite{thong2023skin}. שתי העבודות מצביעות על מקורות מידע שלא נוצלו לכיול צבע במסחר: העין כייחוס, ומרחב צבע רב-ממדי כמטרה.

\begin{keymessagebox}
הספרות פתרה את קביעות הצבע לתמונות כלליות ולתמונות עם ייחוס חיצוני. היא לא פתרה את המקרה שבו העור עצמו הוא רוב הסצנה ואין ייחוס חיצוני. במקרה הזה העין היא הייחוס היחיד שזמין תמיד.
\end{keymessagebox}

\section{תחום ב: שחזור גוף ויד תלת-ממדי}
\label{sec:lit-body}

מודל הגוף \en{SMPL-X} מייצג גוף אנושי, ידיים ופנים כרשת משולשים שנשלטת בשלושה וקטורים: צורה $\boldsymbol{\beta}$, תנוחה $\boldsymbol{\theta}$ והבעה $\boldsymbol{\psi}$ \cite{pavlakos2019smplx}:
\begin{equation}
M(\boldsymbol{\beta},\boldsymbol{\theta},\boldsymbol{\psi}) = W\big(T_P(\boldsymbol{\beta},\boldsymbol{\theta},\boldsymbol{\psi}),\,J(\boldsymbol{\beta}),\,\boldsymbol{\theta},\,\mathcal{W}\big),
\label{eq:smplx}
\end{equation}
כאשר $T_P$ היא תבנית הגוף לאחר עיוותי צורה ותנוחה, $J$ מיקומי המפרקים ו-$W$ פונקציית עיוות לינארית לפי משקלי $\mathcal{W}$. מכיוון שווקטור הצורה $\boldsymbol{\beta}$ מוגדר במרחב מטרי, אפשר לחלץ ממנו היקפים ואורכים. עם זאת, השחזור מתמונה יחידה סובל מעמימות קנה מידה מובנית: גוף גבוה ורחוק נראה כמו גוף נמוך וקרוב.

שיטות הרגרסיה העדכניות, \en{HMR 2.0} \cite{goel2023hmr2} ו-\en{FrankMocap} \cite{rong2021frankmocap}, מייעלות את שגיאת מיקום המפרקים, שאינה מדד למידות: מודל שמצטיין בה עלול לטעות בהיקף המותן בכמה סנטימטרים, ולכן \en{REQ-BOD-03} מוגדרת בהיקפים ונבדקת מול מדידת סרט. רואיס ועמיתיו פונים ישירות לבעיית המדידה: הם אומדים היקפים מצלליות, והגובה של המשתמש משמש אצלם קלט \cite{ruiz2022bodym}. זו ההנחה ש-\en{AURAM} מאמצת: נתון טבלאי אחד, הגובה המדווח, פותר את עמימות קנה המידה. בכף היד, \en{MediaPipe Hands} מאתרת \num{21} נקודות ציון בזמן אמת במכשיר \cite{zhang2020mediapipe} ו-\en{Segment Anything} מפלחת אובייקט לפי הנחיה \cite{kirillov2023sam}; השילוב ביניהן הוא הבסיס ל-\en{REQ-HND-04}.

\section{תחום ג: הדמיה גנרטיבית של לבוש, איפור וציפורניים}
\label{sec:lit-gen}

מודלי דיפוזיה לומדים להפוך תהליך הדרגתי של הוספת רעש לתמונה \cite{ho2020ddpm}. מודלי דיפוזיה סמויה (\en{LDM}) עושים זאת במרחב דחוס של מקודד אוטומטי, ומקטינים את עלות החישוב בסדר גודל \cite{rombach2022ldm}; \en{ControlNet} מוסיפה התניה במבנה מרחבי, כמו שלד תנוחה, בלי לאמן מחדש את המודל הבסיסי \cite{zhang2023controlnet}.

הדמיית לבוש וירטואלית (\en{Virtual Try-On}, \en{VTON}) עברה שני דורות. הראשון עיוות (\en{Warping}) את תמונת הבגד אל הגוף: \en{VITON-HD} ברזולוציה גבוהה \cite{choi2021vitonhd}, ו-\en{GP-VTON} לבגדים מורכבים \cite{xie2023gpvton}. השני מבוסס על דיפוזיה: \en{LaDI-VTON} \cite{morelli2023ladivton}, \en{StableVITON} \cite{kim2024stableviton}, \en{TryOnDiffusion} \cite{zhu2023tryondiffusion}, ו-\en{IDM-VTON}, שמשלבת מודל דיפוזיה גדול עם רשת ייעודית לפרטי הבגד ומשפרת נאמנות בתמונות "מהשטח" \cite{choi2024idmvton}. באיפור, \en{PSGAN} העבירה סגנון בין פנים בתנוחות ובהבעות שונות \cite{jiang2020psgan}.

ההערכה בתחום נשענת על שלושה מדדים: \en{FID} לאיכות התפלגותית \cite{heusel2017fid}, \en{LPIPS} לדמיון תפיסתי לתמונת ייחוס \cite{zhang2018lpips}, ודמיון \en{ArcFace} לשמירת זהות \cite{deng2019arcface}. אף אחד מהם אינו מודד אם צבע הבגד או צבע העור נשמרו באופן מכויל, ואף אחד אינו מודד אם הבגד מוצג במידה הנכונה. ייצוגים נפחיים חדשים, כמו \en{3D Gaussian Splatting} \cite{kerbl2023gaussian}, ותיאור התאורה בהרמוניות כדוריות \cite{ramamoorthi2001sh}, פותחים דור שלישי: בגד שמוצג בתלת-ממד ובתאורה שהמשתמשת בוחרת.

\begin{keymessagebox}
מודלי ההדמיה של הדור השני מצטיינים בפוטו-ריאליזם, אך הם "מלבישים" את הבגד על הצללית שבתמונה. בגד במידה \en{S} ובגד במידה \en{L} ייראו אותו דבר, ושגיאת הצבע של תמונת הקלט עוברת כפי שהיא אל תמונת הפלט.
\end{keymessagebox}

\section{תחום ד: אחזור רב-מודאלי, \en{RAG} וכריית טרנדים}
\label{sec:lit-rag}

\en{CLIP} למד מרחב שיכון משותף לתמונות ולטקסט \cite{radford2021clip}, ו-\en{BLIP-2} חיבר מקודד תמונה קפוא למודל שפה גדול \cite{li2023blip2}; יחד הם מאפשרים לתאר פריט במילים ולמצוא אותו בקטלוג. יצירה מועשרת-אחזור (\en{RAG}) מעגנת את התשובה במאגר חיצוני עדכני \cite{lewis2020rag}, ו-\en{FAISS} מחפשת שכנים קרובים בין מיליוני פריטים באלפיות שנייה \cite{johnson2021faiss}. בצד התוכן, רוב מערכות ההמלצה לאופנה נשענות על אינטראקציות עבר, ומעט מהן משלבות את מאפייני הגוף של הלקוחה \cite{deldjoo2023fashionrec}. חיזוי טרנדים נחקר כבעיית סדרות עתיות \cite{alhalah2017fashionforward,ma2020fashiontrend}, אך ברמת האוכלוסייה: אף עבודה אינה שואלת אם הטרנד מתאים לגוון העור או למבנה הגוף של לקוחה מסוימת.

\section{שיטה, תרומה ומגבלה: סיכום ההשוואה}
\label{sec:lit-summary}

טבלה~\ref{tab:litsum} מרכזת לכל מקור מרכזי את השיטה, התרומה והמגבלה ביחס לבעיה של \en{AURAM}. עמודת המגבלה היא החשובה: היא מתעדת את ההנחות המפשטות שהפרויקט נדרש להסיר, והן חוזרות על שלושה דפוסים (סעיף~\ref{sec:synthesis}).

\begin{reporttable}[\footnotesize]{שיטה, תרומה עיקרית ומגבלה ביחס לבעיה, לפי מקור}{tab:litsum}{P{2.9cm} R R R}
\headrow \hc{מקור} & \hc{שיטה} & \hc{תרומה עיקרית} & \hc{מגבלה ביחס לבעיה} \\ \midrule
\en{FFCC} \cite{barron2017ffcc} & היסטוגרמה לוג-כרומטית ופורייה & אמידת מאיר מהירה ומדויקת & מניחה סצנה מגוונת, לא סלפי \\
תיקון איזון לבן \cite{afifi2019wb} & תיקון לא-לינארי מבוסס דוגמאות & מתקן תמונות שאוזנו בשגיאה & אינו מנצל עוגן בגוף \\
העין כגשש אור \cite{nishino2004eyes} & מודל גיאומטרי של הקרנית & שחזור תאורת הסביבה מהעין & לא יושם לכיול צבע עור \\
גוון עור רב-ממדי \cite{thong2023skin} & מדד בהירות וגוון & הערכת הוגנות לפי גוון & מדד הערכה, לא שיטת כיול \\
\en{SMPL-X} \cite{pavlakos2019smplx} & מודל גוף פרמטרי & גוף, ידיים ופנים במודל אחד & עמימות קנה מידה מתמונה \\
מדידת גוף \cite{ruiz2022bodym} & רגרסיה מצלליות וגובה & היקפים מתמונות & מחייב צללית נקייה \\
\en{MediaPipe Hands} \cite{zhang2020mediapipe} & רשת קלה לנקודות ציון & איתור יד בזמן אמת בנייד & ללא פילוח ציפורניים \\
\en{ControlNet} \cite{zhang2023controlnet} & התניה מרחבית בדיפוזיה & שליטה במבנה בלי אימון מחדש & אינה מבטיחה נאמנות צבע \\
\en{StableVITON} \cite{kim2024stableviton} & התאמה סמנטית בדיפוזיה סמויה & שימור פרטי בגד & סטודיו, ללא מידות \\
\en{IDM-VTON} \cite{choi2024idmvton} & דיפוזיה עם רשת בגד מקבילה & נאמנות בתמונות מהשטח & ללא כיול, ללא מידות, יקרה \\
\en{CLIP} \cite{radford2021clip} & שיכון משותף בלמידה ניגודית & אחזור תמונה-טקסט כללי & אינו מבחין בתת-גוון \\
\en{RAG} \cite{lewis2020rag} & אחזור ויצירה משולבים & תשובות מעוגנות במקור & נבחן על טקסט, לא על אופנה \\
\en{Fashion Forward} \cite{alhalah2017fashionforward} & סדרות עתיות של סגנונות & חיזוי פופולריות & ברמת אוכלוסייה בלבד \\
\end{reporttable}

\csname @nobreaktrue\endcsname % one-off (v3.5): the cheap \needspace break left 35-40% of the page blank
\section{מטריצת הספרות}
\label{sec:matrix}

טבלה~\ref{tab:litmatrix} מרכזת את הסקירה. השורות הן המקורות והעמודות הן שש יכולות שהפרויקט דורש. עיגול מלא מסמן יכולת שנבחנה במלואה, עיגול ריק יכולת חלקית, ומשבצת ריקה יכולת שלא נבחנה.

\begin{reporttable}[\footnotesize]{מטריצת הספרות: מקורות מול שש יכולות. המשבצות הריקות הן הפער}{tab:litmatrix}{R Q{1.45cm} Q{1.45cm} Q{1.45cm} Q{1.45cm} Q{1.45cm} Q{1.45cm}}
\headrow \hc{מקור} & \hc{כיול צבע} & \hc{עוגן ביומטרי} & \hc{גוף מטרי} & \hc{הדמיה} & \hc{אחזור וטרנד} & \hc{עלות ונייד} \\ \midrule
\en{FFCC} \cite{barron2017ffcc} & \cellok & \cellno & \cellno & \cellno & \cellno & \cellpartial \\
\en{FC4} \cite{hu2017fc4} & \cellok & \cellno & \cellno & \cellno & \cellno & \cellno \\
תיקון איזון לבן \cite{afifi2019wb} & \cellok & \cellno & \cellno & \cellno & \cellno & \cellno \\
פולינום-שורש \cite{finlayson2015rootpoly} & \cellok & \cellno & \cellno & \cellno & \cellno & \cellpartial \\
העין כגשש אור \cite{nishino2004eyes} & \cellpartial & \cellok & \cellno & \cellno & \cellno & \cellno \\
גוון עור רב-ממדי \cite{thong2023skin} & \cellpartial & \cellpartial & \cellno & \cellno & \cellno & \cellno \\
\en{SMPL-X} \cite{pavlakos2019smplx} & \cellno & \cellno & \cellpartial & \cellno & \cellno & \cellno \\
\en{HMR 2.0} \cite{goel2023hmr2} & \cellno & \cellno & \cellpartial & \cellno & \cellno & \cellno \\
מדידת גוף \cite{ruiz2022bodym} & \cellno & \cellno & \cellok & \cellno & \cellno & \cellpartial \\
\en{MediaPipe Hands} \cite{zhang2020mediapipe} & \cellno & \cellno & \cellpartial & \cellno & \cellno & \cellok \\
\en{PSGAN} \cite{jiang2020psgan} & \cellno & \cellno & \cellno & \cellpartial & \cellno & \cellno \\
\en{GP-VTON} \cite{xie2023gpvton} & \cellno & \cellno & \cellno & \cellok & \cellno & \cellno \\
\en{StableVITON} \cite{kim2024stableviton} & \cellno & \cellno & \cellno & \cellok & \cellno & \cellno \\
\en{TryOnDiffusion} \cite{zhu2023tryondiffusion} & \cellno & \cellno & \cellpartial & \cellok & \cellno & \cellno \\
\en{IDM-VTON} (דגל) \cite{choi2024idmvton} & \cellno & \cellno & \cellno & \cellok & \cellno & \cellpartial \\
\en{CLIP} \cite{radford2021clip} & \cellno & \cellno & \cellno & \cellno & \cellok & \cellno \\
\en{RAG} \cite{lewis2020rag} & \cellno & \cellno & \cellno & \cellno & \cellok & \cellno \\
\en{Fashion Forward} \cite{alhalah2017fashionforward} & \cellno & \cellno & \cellno & \cellno & \cellok & \cellno \\
סקירת המלצות אופנה \cite{deldjoo2023fashionrec} & \cellno & \cellno & \cellpartial & \cellno & \cellok & \cellno \\
\hlrow \textbf{הפרויקט המוצע, \en{AURAM}, סוגר את הפער} & \cellok & \cellok & \cellok & \cellok & \cellok & \cellok \\
\end{reporttable}

הקריאה לאורך העמודות מלמדת יותר מהקריאה לאורך השורות: אף מקור אינו מכסה יותר משתי יכולות במלואן, ואף מקור בתחום ההדמיה אינו מכסה את כיול הצבע או את העוגן הביומטרי. עמודת העלות והפריסה בנייד ריקה כמעט לגמרי בתחום ההדמיה, כי המאמרים מדווחים על איכות ולא על עלות ליחידה. השורה האחרונה היא התחייבות ולא הישג.

\csname @nobreaktrue\endcsname % one-off (v3.5): the cheap \needspace break left 35-40% of the page blank
\section{סינתזה ביקורתית: שלוש הנחות שחוזרות בספרות}
\label{sec:synthesis}

מעבר למשבצות הריקות, הסקירה חושפת שלוש הנחות מפשטות שחוזרות בכל התחומים. הן מסבירות למה המשבצות ריקות.

\textbf{הנחה \num{1}: הקלט כבר מכויל.} מודלי ההדמיה, מנועי ההמלצה ומודלי הטרנדים מקבלים את צבעי התמונה כנתון. רק ספרות קביעות הצבע מטפלת בהם, והיא אינה מחוברת לשאר הצינור. בגלל ההנחה הזו, שגיאה של \num{$\Delta E_{00} \approx 5$} בקלט עוברת בלי מגע עד לפלט.

\textbf{הנחה \num{2}: צורה במקום מידה.} מודלי הגוף מייעלים דיוק של תנוחה, ומודלי ההדמיה מייעלים התאמה לצללית. שניהם עובדים ביחידות של פיקסלים, ואף אחד אינו מחויב לסנטימטרים. לכן גם מודל מצוין אינו יכול לומר אם בגד במידה \en{M} יתאים.

\textbf{הנחה \num{3}: עלות החישוב אינה חלק מהבעיה.} המאמרים בתחום ההדמיה מדווחים על איכות ולא על עלות ליחידה, ואינם מגדירים תקציב חישוב. זו הנחה סבירה במחקר, אך במוצר היא קובעת אם הרכיב יופעל בכלל.

לצד שלוש ההנחות קיימת גם ביקורת מתודולוגית. רוב מאגרי ההדמיה צולמו בסטודיו, עם רקע אחיד ותאורה מבוקרת, ולכן אינם מייצגים את התנאים של משתמשת בביתה. בנוסף, עבודות על גוון עור הראו שמדדים שנשענים על בהירות בלבד מסתירים הבדלים בין קבוצות \cite{thong2023skin}. מכאן נגזרים שני עקרונות לתכנון הניסוי בפרויקט: איסוף בתנאי בית אמיתיים בארבע תאורות, ודיווח של תוצאות לפי תת-קבוצות גוון ולא רק בממוצע.

\csname @nobreaktrue\endcsname % one-off: the cheap \needspace break left 40% of the page blank
\section{בחירת מאמר הדגל}
\label{sec:flagship}

מאמר הדגל (\en{Flagship Paper}) הוא הפרסום שממנו הפרויקט שואב את הבסיס המתודולוגי הישיר ביותר. הוא נבחר בחיתוך של שלושה קריטריונים (איור~\ref{fig:venn}): במה שפיטה בדירוג גבוה, קרבה ישירה לבעיה, וחזרתיות (\en{Replicability}) שמאפשרת לשחזר את התוצאות \cite{segal2026booklet}.

\begin{reportfigure}{מאמר הדגל נמצא בחיתוך של שלושת הקריטריונים}{fig:venn}
\begin{tikzpicture}[x=1cm,y=1cm,scale=0.76,
  cat/.style={fill=#1,text=brandcream,font=\small\bfseries,rounded corners=7pt,inner xsep=7pt,inner ysep=3pt},
  pap/.style={fill=white,draw=#1,line width=0.6pt,font=\scriptsize,rounded corners=2pt,inner sep=2pt,align=center}]
  % three distinct hues: Velvet Forest (venue), Copper (closeness), Dust (replicability)
  \fill[brandprimary,fill opacity=0.16] (0,1.1) circle (2.6);
  \fill[brandhighlight,fill opacity=0.20] (-1.4,-1.0) circle (2.6);
  \fill[brandsecondary,fill opacity=0.22] (1.4,-1.0) circle (2.6);
  \draw[brandprimary,line width=1.4pt] (0,1.1) circle (2.6);
  \draw[brandhighlight,line width=1.4pt] (-1.4,-1.0) circle (2.6);
  \draw[brandsecondary,line width=1.4pt] (1.4,-1.0) circle (2.6);
  % criterion names outside the circles
  \node[cat=brandprimary,anchor=south] at (0,3.78) {\heb{במה שפיטה}};
  \node[cat=brandhighlight,anchor=north] at (-2.9,-3.68) {\heb{קרבה לבעיה}};
  \node[cat=brandsecondary,anchor=north] at (2.9,-3.68) {\heb{חזרתיות}};
  % papers, each label fully inside its own region
  \node[pap=brandprimary] at (-1.65,0.8) {\en{TryOn}\\\en{Diffusion}};
  \node[pap=brandsecondary] at (1.6,0.9) {\en{ControlNet}\\\en{GP-VTON}};
  \node[pap=brandsecondary] at (0,-1.95) {\en{VITON-HD}};
  \node[fill=brandhighlight,text=white,font=\footnotesize\bfseries,rounded corners=3pt,inner sep=2.5pt] at (0,-0.8) {\en{IDM-VTON}};
\end{tikzpicture}
\end{reportfigure}

טבלה~\ref{tab:flagcand} מציגה את ההכרעה: \en{TryOnDiffusion} עומדת בשני הקריטריונים הראשונים, אך המודל והנתונים שלה אינם זמינים לציבור ואי אפשר לשחזר את תוצאותיה; רק \en{IDM-VTON} עומדת בשלושתם.

\begin{reporttable}{המועמדים למאמר הדגל מול שלושת הקריטריונים}{tab:flagcand}{R Q{2.8cm} Q{3.9cm} Q{3.2cm}}
\headrow \hc{מועמד} & \hc{במה שפיטה} & \hc{קרבה לבעיה} & \hc{חזרתיות} \\ \midrule
\en{TryOnDiffusion} \cite{zhu2023tryondiffusion} & \en{CVPR 2023} \cellok & גבוהה \cellok & קוד ונתונים סגורים \cellno \\
\en{StableVITON} \cite{kim2024stableviton} & \en{CVPR 2024} \cellok & גבוהה \cellok & קוד פתוח \cellok \\
\en{GP-VTON} \cite{xie2023gpvton} & \en{CVPR 2023} \cellok & בינונית, דור העיוות \cellpartial & קוד פתוח \cellok \\
\hlrow \en{IDM-VTON} \cite{choi2024idmvton} & \en{ECCV 2024} \cellok & גבוהה, תמונות מהשטח \cellok & קוד ומשקלים \cellok \\
\end{reporttable}

ההכרעה בין \en{StableVITON} ל-\en{IDM-VTON} נשענת על הקרבה לבעיה: \en{IDM-VTON} מתמקדת במפורש בתמונות שצולמו מחוץ לסטודיו, וזה המצב של המשתמשת ב-\en{AURAM}. הבחירה אינה סופית: אם ניסוי השחזור בשבוע \num{3} ייכשל, \en{StableVITON} היא החלופה המוגדרת מראש.

\section{ניתוח מאמר הדגל והגדרת פער הידע}
\label{sec:flagship-analysis}

\subsection{המודל והמשוואות השולטות}
\en{IDM-VTON} בנויה על מודל דיפוזיה סמויה גדול לציור-מחדש (\en{Inpainting}), שמקבל כקלט את תמונת האדם שבה אזור הבגד ממוסך, מפת תנוחה, ואת תמונת הבגד \cite{choi2024idmvton}. האימון ממזער את שגיאת חיזוי הרעש במרחב הסמוי:
\begin{equation}
\mathcal{L} = \mathbb{E}_{\mathbf{z}_0,\boldsymbol{\epsilon},t}\Big[\big\lVert \boldsymbol{\epsilon} - \boldsymbol{\epsilon}_\theta(\mathbf{z}_t,\,t,\,\mathbf{x}_m,\,\mathbf{p},\,\mathbf{g}) \big\rVert_2^2\Big],
\label{eq:ldmloss}
\end{equation}
כאשר $\mathcal{L}$ הוא ההפסד הממוזער, $\boldsymbol{\epsilon}_\theta$ רשת חיזוי הרעש עם הפרמטרים $\theta$, $\mathbf{z}_t$ הייצוג הסמוי הרועש, $\mathbf{x}_m$ תמונת האדם הממוסכת, $\mathbf{p}$ התנוחה ו-$\mathbf{g}$ הבגד. הבגד נכנס למודל בשני מסלולים. מסלול סמנטי ברמה גבוהה עובר דרך מתאם הנחיה חזותית, שמוסיף לתשומת הלב הצולבת של הטקסט רכיב מקביל של תמונה:
\begin{equation}
\mathbf{Z} = \operatorname{Attn}(\mathbf{Q},\mathbf{K}_{\text{txt}},\mathbf{V}_{\text{txt}}) + \lambda\,\operatorname{Attn}(\mathbf{Q},\mathbf{K}_{\text{img}},\mathbf{V}_{\text{img}}).
\label{eq:ipadapter}
\end{equation}
כאשר $\mathbf{Q}$ הן השאילתות ממאפייני התמונה הרועשת, $\mathbf{K}$ ו-$\mathbf{V}$ המפתחות והערכים מהטקסט או מתמונת הבגד, $\lambda$ משקל המסלול החזותי, ו-$\mathbf{Z}$ הפלט המשולב. מסלול שני, ברמה נמוכה, עובר דרך רשת \en{UNet} מקבילה (\en{GarmentNet}) שמחלצת מאפיינים מפורטים של הבגד, כמו לוגו, טקסטורה ותבנית. מאפיינים אלה מצורפים למאפייני האדם בתוך שכבות תשומת הלב העצמית, כך שכל פיקסל של הגוף יכול "להסתכל" על כל פיקסל של הבגד. בזמן ההסקה נעשה שימוש בהכוונה ללא מסווג (\en{Classifier-Free Guidance}):
\begin{equation}
\tilde{\boldsymbol{\epsilon}}_\theta = \boldsymbol{\epsilon}_\theta(\mathbf{z}_t,\varnothing) + w\,\big(\boldsymbol{\epsilon}_\theta(\mathbf{z}_t,\mathbf{c}) - \boldsymbol{\epsilon}_\theta(\mathbf{z}_t,\varnothing)\big),
\label{eq:cfg}
\end{equation}
כאשר $\mathbf{c}$ הוא תנאי ההתניה, $\varnothing$ התניה ריקה ו-$w$ עוצמת ההכוונה. השיטה מכפילה את מספר הערכות הרשת בכל צעד. איור~\ref{fig:flagship} מסכם את הצינור ומסמן את שלוש נקודות התורפה.

\begin{reportfigure}{צינור מאמר הדגל \en{IDM-VTON} ושלוש מגבלותיו ביחס לבעיה}{fig:flagship}
\begin{tikzpicture}[x=1cm,y=1cm,
  m/.style={draw=brandprimary,fill=brandtint,rounded corners=3pt,minimum height=0.85cm,align=center,font=\scriptsize,text width=2.2cm},
  l/.style={draw=brandwarn,fill=brandwarn!8,rounded corners=2pt,font=\scriptsize,align=center,text width=2.6cm}]
  \node[m] (in) at (0,0) {\heb{תמונת האדם}\\\heb{(צבע לא מכויל)}};
  \node[m] (g) at (0,-1.35) {\heb{תמונת הבגד}};
  \node[m] (gn) at (3.0,-1.35) {\en{GarmentNet}\\\en{+ IP-Adapter}};
  \node[m] (un) at (3.0,0) {\en{SDXL Inpainting}\\\en{TryonNet}};
  \node[m] (out) at (6.1,0) {\heb{תמונת הדמיה}};
  \node[l] (l1) at (0,1.35) {\heb{1: אין כיול צבע}};
  \node[l] (l2) at (6.1,1.35) {\heb{2: אין מידות מטריות}};
  \node[l] (l3) at (6.1,-1.35) {\heb{3: עלות ושניות לתמונה}};
  \draw[-Stealth,thick] (in) -- (un); \draw[-Stealth,thick] (g) -- (gn);
  \draw[-Stealth,thick] (gn) -- (un); \draw[-Stealth,thick] (un) -- (out);
  \draw[brandwarn,dashed] (l1) -- (in); \draw[brandwarn,dashed] (l2) -- (out); \draw[brandwarn,dashed] (l3) -- (un);
\end{tikzpicture}
\end{reportfigure}

\subsection{סביבת הניסוי ומגבלות ההנחות}
המאמר מעריך את המודל על מאגרי \en{VITON-HD} ו-\en{DressCode}, שבהם הדוגמניות מצולמות בסטודיו ברקע אחיד ובתאורה מבוקרת, ועל אוסף של תמונות "מהשטח" \cite{choi2024idmvton}. מהניתוח עולות שלוש מגבלות מהותיות ביחס לבעיה שהוגדרה בפרק~\ref{ch:problem}, והן מסומנות באיור~\ref{fig:flagship}:

\textbf{מגבלה \num{1}, הנחת תאורה.} המודל מעתיק את צבעי תמונת הקלט: בתאורת להט העור והבגד יוצגו בגוון חמים, והמדדים של המאמר (\en{FID}, \en{LPIPS}) אינם רגישים לכך, כי גם תמונת הייחוס צולמה באותה תאורה. \textbf{מגבלה \num{2}, היעדר מידות.} הבגד מותאם לצללית שבתמונה, בלי מידת הבגד ובלי היקפי הגוף; ההדמיה עונה על השאלה "איך זה ייראה" ולא על השאלה "האם זה יתאים". \textbf{מגבלה \num{3}, משאבי חישוב.} לפי אומדן הצוות, שיימדד בניסוי השחזור, ההסקה נמשכת יותר מ-\num{8 s} לתמונה גם על מעבד גרפי מדרגת \en{A100}, רחוק מ-\en{REQ-GEN-06} ומ-\en{REQ-BIZ-12}.

\subsection{מה הפרויקט לוקח, משנה ומוסיף}
טבלה~\ref{tab:takeadd} מפרטת מה נלקח ממאמר הדגל כמו שהוא, מה משתנה ומה מתווסף: הניתוח מוביל לשימוש מבוקר בו, ולא לדחייתו.

\begin{reporttable}{היחס למאמר הדגל: מה נלקח, מה משתנה ומה מתווסף}{tab:takeadd}{P{1.8cm} R R}
\headrow \hc{פעולה} & \hc{רכיב} & \hc{נימוק} \\ \midrule
לוקחים & רשת הבגד המקבילה ומנגנון תשומת הלב & הנאמנות לפרטי הבגד היא החוזקה המוכחת \\
לוקחים & הקוד והמשקלים שפורסמו & מאפשרים שחזור בשבוע \num{3} \\
משנים & מספר צעדי הדגימה: זיקוק ל-\num{8} & עמידה ב-\en{REQ-GEN-06} וב-\en{REQ-BIZ-12} \\
משנים & תמונת הקלט מכוילת לפני ההדמיה & סגירת מגבלה \num{1} \\
מוסיפים & התניה במידה ובמפת היקפים & סגירת מגבלה \num{2} \\
מוסיפים & מסיכת יד וציפורניים כהתניה & הפחתת \en{RSK-GEN-02} \\
\end{reporttable}

הפרויקט אינו מתחרה במאמר הדגל, אלא עוטף אותו. הנאמנות הגנרטיבית נשארת של המודל, והפרויקט מוסיף את מה שחסר לו כדי לשמש במוצר: קלט מכויל, מידה ומסגרת עלות.

\subsection{פער הידע}
שלוש המגבלות, יחד עם המשבצות הריקות במטריצה, מתנקזות למשפט אחד:

\begin{keymessagebox}[פער הידע]
הפתרונות הקיימים להדמיה ולהמלצה מניחים שצבעי תמונת הסלפי נכונים, ובמקורות שנסקרו לא נמצאה עבודה שבודקת אם העין, בעזרת לובן העין וצבע עיניים מדווח, יכולה לשמש עוגן לכיול גוון העור בתאורה לא מבוקרת.
\end{keymessagebox}

הפער נבנה בדרך של ערעור על הנחה (\en{Problematization}) ולא רק באיתור משבצת ריקה, במובן שהציעו סנדברג ואלבסון \cite{sandberg2011gap}: ההנחה שצבעי תמונת הקלט נכונים משותפת לכל הספרות של תחום ג, ולכן שיפור של מודל ההדמיה לבדו לא יסגור אותו. ההשערות בסעיף~\ref{sec:rq} בוחנות את הפער הזה ישירות. מידות הגוף והמימון אינם חלק מטענת הפער אלא רכיבי אינטגרציה. הטענה מוגבלת למקורות שנסקרו, וחיפוש שיטתי (\en{IEEE Xplore}, \en{ACM DL}, \en{Google Scholar}) מתוכנן בחבילת העבודה הראשונה. הפרויקט סוגר את הפער בשכבת הכיול, ומעביר למודל הדגל קלט מכויל ומידות מטריות.

\subsection{הרצה מקדימה של מבחני ההשערות}
\label{sec:pilot}
כדי לוודא שתכנון הניתוח בסעיף~\ref{sec:rq} ישים, המודל במשוואה~\eqref{eq:hyp} הותאם לניסוי מדומה בגודל המתוכנן: \num{120} משתתפות, כל אחת בארבע משפחות התאורה, מהסימולציה הספקטרלית (\enlong{scripts/statistical\_validation.py}). אומדן השיפור הוא \num{$\hat\beta_1 = -4.43$}, עם רווח סמך של \num{95\%} \num{[-4.68, -4.18]}. רכיבי השונות הם \num{0.79} למשתתפת, \num{1.92} לתאורה ו-\num{3.81} לשארית, כלומר מתאם תוך-משתתפת של \num{0.12}, נמוך מטווח התכנון, ולכן העוצמה ל-$H_{0a}$ רק גדלה; ל-$H_{0b}$ העוצמה נמוכה כל עוד החציון קרוב ל-\num{1.8}. סטיית התקן של השארית גבוהה יותר בעור בהיר (\num{2.07}) מאשר בעור כהה (\num{1.69}), ולכן ההנחה של שונות שארית שונה בין קבוצות הגוון מוצדקת. $H_{0a}$ נדחית. $H_{0b}$ אינה נדחית: חציון השגיאה הוא \num{1.94}, והגבול העליון שלו, \num{2.06}, מעל \num{1.8}. ערכי $p$ על נתונים סינתטיים אינם ראיה ואינם מדווחים; ההרצה מאמתת את צינור הניתוח ומראה שיעד החציון אינו מובטח (\en{RSK-ALG-07}).

\begin{examplebox}[בקול שלנו: עור אינו מטריצת \en{RGB}]
מהנדסי ראייה ממוחשבת נוטים לראות בעור מטריצה של שלושה ערוצים שצריך "לתקן". אנחנו רואות בו רקמה ביולוגית וזהות תרבותית: גוון שעליו אישה בוחרת שפתון, ונער בוחר חולצה. מכאן גם הדרך שלנו לקרוא תוצאה: אי-דחייה של $H_{0b}$ אינה כישלון אלא מידע, והיא אומרת לנו בדיוק איפה להשקיע את שבועות המימוש.
\end{examplebox}
